%==========================================================================================
\section{\texttt{xrf\_beamline}: the incident spectrum}\label{sec:beamline}
%==========================================================================================

\subsection{Overview}

Under white-beam (Laue) illumination every element is excited by the whole incident
spectrum $\Phi(E)$, and the relative intensities of its fluorescence lines---not only its
total signal---depend on the shape of that spectrum: each L or M subshell is ionised only by
photons above its own edge, and those edges can be kilo-electronvolts apart (Ir: L$_3$ at
11.22\,keV, L$_1$ at 13.42\,keV). The module computes $\Phi(E)$ at the sample from first
principles: the synchrotron radiation of a bending magnet \citep{Schwinger1949,Kim1989},
accepted through the apertures of the beamline, and filtered by grazing-incidence mirrors
\citep{Parratt1954,Nevot1980}, beryllium windows and the air path. It also provides narrow
``monochromatic'' spectra for tube or monochromator excitation, and the air model used by
the detector-efficiency calculation of \fn{xrf_fit}.

The quantities returned by \fn{Beamline.spectrum} are fluxes per mrad of horizontal angle
and per 0.1\,\% bandwidth; only their \emph{shape} matters for line ratios and for the
matrix corrections, but the absolute scale is correct (photons/s) for a given ring current.

\begin{hyp}
\begin{itemize}[nosep,leftmargin=1em]
\item Bending-magnet radiation of a filament electron beam (no emittance, no energy spread);
  the horizontal fan is uniform, so horizontal apertures scale the flux but not the spectrum.
\item Mirrors: specular s-polarised Fresnel reflectivity of a homogeneous coating of nominal
  density, with an optional N\'evot--Croce roughness factor; no slope errors, no thermal
  deformation, no contamination layer. Bent mirrors act as thin cylindrical lenses in the
  paraxial approximation.
\item Windows and air: Beer--Lambert attenuation with total cross-sections; dry air, ideal
  gas.
\item The BM32 defaults (magnetic field 0.86\,T, Ir coatings, $2\times500$\,\textmu m Be,
  3.3\,m of air, element positions) were confirmed in October 2026; mirror roughness is
  unknown (0) and the KB grazing angles are not recorded in the scan files (3\,mrad).
\end{itemize}
\end{hyp}

\pycode{xrf_beamline}{_header}

\subsection{Source}

\api[\texttt{BendingMagnetSource}]{\texttt{BendingMagnetSource(ring\_energy\_gev=6.04, magnetic\_field\_t=0.86, current\_a=0.2)}}
Synchrotron radiation of a bending magnet. The defaults describe the BM32 source on ESRF-EBS.
\begin{args}
\item[ring\_energy\_gev] electron energy $E_e$ (6.04\,GeV for EBS).
\item[magnetic\_field\_t] dipole field $B$; sets the critical energy, to which the high-energy
  end of the spectrum is exponentially sensitive.
\item[current\_a] ring current $I$, in amperes.
\end{args}
\blk{Physics} The critical energy and the Lorentz factor are
\begin{equation}
  E_c\,[\mathrm{keV}] = 0.665\,E_e^2[\mathrm{GeV}]\,B[\mathrm{T}], \qquad
  \gamma = \frac{E_e}{m_e c^2} \approx 11\,820 ,
\end{equation}
giving $E_c = 20.9$\,keV for BM32. With $y = E/E_c$, the flux integrated over the vertical
opening angle, per mrad of horizontal angle and per 0.1\,\% bandwidth, is \citep{Kim1989}
\begin{equation}
  \frac{\dd N}{\dd\theta} = 2.457\times10^{13}\,E_e[\mathrm{GeV}]\,I[\mathrm{A}]\;G_1(y),
  \qquad G_1(y) = y\int_y^\infty K_{5/3}(x)\,\dd x ,
  \label{eq:g1}
\end{equation}
and the angular density at vertical angle $\psi$, per mrad$^2$, with $X = \gamma\psi$ and
$\xi = \tfrac{y}{2}(1+X^2)^{3/2}$, is
\begin{equation}
  \frac{\dd^2 N}{\dd\theta\,\dd\psi} = 1.327\times10^{13}\,E_e^2\,I\;y^2(1+X^2)^2
  \left[K_{2/3}^2(\xi) + \frac{X^2}{1+X^2}K_{1/3}^2(\xi)\right],
  \label{eq:angular}
\end{equation}
where the first term is the horizontally ($\sigma$) polarised part and the second the
vertically ($\pi$) polarised part. On axis it reduces to $1.327\times10^{13}E_e^2 I\,H_2(y)$
with $H_2(y) = y^2K_{2/3}^2(y/2)$. The vertical opening angle is of order $1/\gamma =
85$\,\textmu rad, and wider for soft photons, so an aperture that accepts only a few
\textmu rad selects the harder on-axis spectrum.
\pycode{xrf_beamline}{BendingMagnetSource}

\api[\texttt{critical\_energy\_kev}, \texttt{gamma} ; \texttt{\_g1}]{\texttt{critical\_energy\_kev}, \texttt{gamma} (properties); \texttt{\_g1(y)} (static)}
$E_c$ and $\gamma$ as above. \fn{_g1} evaluates $G_1(y)$ by adaptive quadrature of
$K_{5/3}$ (\fn{scipy.special.kv}) from $y$ to infinity; it returns 0 for $y\le0$ and for
$y>50$, where $G_1 < 10^{-20}$.
\pycode{xrf_beamline}{BendingMagnetSource.critical_energy_kev}
\pycode{xrf_beamline}{BendingMagnetSource.gamma}
\pycode{xrf_beamline}{BendingMagnetSource._g1}

\api[\texttt{spectral\_flux}]{\texttt{spectral\_flux(energies\_kev)}}
Eq.~\eqref{eq:g1} on an energy grid (keV): photons/s/mrad(horizontal)/0.1\,\%\,BW, integrated
over the vertical angle. Returns an array of the grid's size.
\pycode{xrf_beamline}{BendingMagnetSource.spectral_flux}

\api[\texttt{angular\_flux}]{\texttt{angular\_flux(energies\_kev, psi\_urad)}}
Eq.~\eqref{eq:angular}: photons/s/mrad$^2$/0.1\,\%\,BW, both polarisations, on a grid of
energies and vertical angles (\textmu rad). Returns an array of shape $(n_E, n_\psi)$.
Bessel functions that overflow or underflow far from the axis are set to zero.
\pycode{xrf_beamline}{BendingMagnetSource.angular_flux}

\api[\texttt{accepted\_flux}]{\texttt{accepted\_flux(energies\_kev, vertical\_acceptance\_urad=None, vertical\_offset\_urad=0.0, n\_angles=41)}}
Flux per mrad (horizontal) per 0.1\,\%\,BW inside the vertical window
$\psi_0 \pm \Delta\psi/2$: the trapezoidal integral of Eq.~\eqref{eq:angular} over
\fn{n_angles} angles, converted from \textmu rad to mrad. \fn{None} (or a non-finite value)
returns the full integral, Eq.~\eqref{eq:g1}. For $\Delta\psi \ll 1/\gamma$ the result is the
on-axis spectrum times $\Delta\psi$.
\pycode{xrf_beamline}{BendingMagnetSource.accepted_flux}

\subsection{Optical elements}

\api[\texttt{Mirror}]{\texttt{Mirror(material, density\_g\_cm3, grazing\_angle\_mrad, roughness\_nm=0.0, label="", distance\_m=None, length\_m=None, plane="vertical", bending\_radius\_m=inf, image\_distance\_m=None, angle\_spread\_mrad=0.0)}}
A grazing-incidence mirror. Its spectral effect is the reflectivity; if it is positioned
(\fn{distance_m} and \fn{length_m}), its footprint also acts as an aperture; if it is bent,
it focuses the beam divergence; if it is a focusing (KB) mirror, its grazing angle varies
along the footprint.
\begin{args}
\item[material, density\_g\_cm3] coating, e.g. \fn{"Ir"}, 22.56.
\item[grazing\_angle\_mrad] angle between the beam and the surface, $\theta$.
\item[roughness\_nm] r.m.s. roughness $\sigma$ (N\'evot--Croce factor; 0 disables it).
\item[distance\_m, length\_m] distance from the source and optical length.
\item[plane] \fn{"vertical"} or \fn{"horizontal"}: deflection plane.
\item[bending\_radius\_m] meridional radius $R$ (\fn{inf} for a flat mirror).
\item[image\_distance\_m] distance $q$ to the focus of a focusing mirror.
\item[angle\_spread\_mrad] half-range of grazing angles over the illuminated footprint;
  filled in by \fn{Beamline.resolved_optics} from \fn{image_distance_m}.
\end{args}
\pycode{xrf_beamline}{Mirror}
\pycode{xrf_beamline}{Mirror.__post_init__}

\api[\texttt{focal\_length\_m}, \texttt{aperture\_m}]{\texttt{focal\_length\_m}, \texttt{aperture\_m} (properties)}
A cylindrically bent mirror of radius $R$ at grazing angle $\theta$ focuses in its
deflection plane with focal length $f = R\sin\theta/2$ (the meridional Coddington equation
at grazing incidence). The footprint of length $L$ accepts a beam of size $L\sin\theta$.
\pycode{xrf_beamline}{Mirror.focal_length_m}
\pycode{xrf_beamline}{Mirror.aperture_m}

\api[\texttt{Mirror.transfer}, \texttt{Mirror.\_reflectivity}]{\texttt{Mirror.transfer(energies\_kev)}, \texttt{Mirror.\_reflectivity(energies\_kev, angle\_mrad)}}
\blk{Physics} With the complex refractive index $n = 1-\delta+i\beta$ of the coating (xraylib
\fn{Refractive_Index_Re} and \fn{_Im}), the s-polarised Fresnel amplitude at grazing angle
$\theta$ is \citep{Parratt1954,AlsNielsen2011}
\begin{equation}
  r_s = \frac{\sin\theta - q}{\sin\theta + q}, \qquad q = \sqrt{n^2-\cos^2\theta},
  \qquad R = |r_s|^2 .
\end{equation}
Total external reflection holds below the critical angle $\theta_c \approx \sqrt{2\delta}$;
since $\delta\propto E^{-2}$, a mirror at fixed angle cuts off above $E_c^\mathrm{mirror}
\propto 1/\theta$ (for Ir, $\delta = 3.4\times10^{-5}$ at 10\,keV, so $\theta_c = 8.3$\,mrad at
10\,keV and the edge of reflection at 3\,mrad lies near 27\,keV). Roughness multiplies the
\emph{amplitude} by $\exp(-2k_{z0}k_{z1}\sigma^2)$ \citep{Nevot1980}, with $k_{z0} =
k_0\sin\theta$ and $k_{z1} = k_0 q$, so the intensity factor is
$\exp(-4k_0^2\sin\theta\,\Re(q)\,\sigma^2)$; below the critical angle $\Re q \ll \sin\theta$
and the factor tends to 1, while well above it the Debye--Waller form
$\exp(-4k_{z0}^2\sigma^2)$ is recovered. Synchrotron radiation is polarised in the orbit
plane and the mirrors of BM32 deflect vertically, so s-polarisation is the relevant case
(at mrad angles both polarisations give the same result anyway).

\fn{transfer} returns $R(E)$, averaged uniformly over 21 angles in $\theta \pm$
\fn{angle_spread_mrad} if a spread is set. Energies where xraylib has no data get $n=1$.
\pycode{xrf_beamline}{Mirror.transfer}
\pycode{xrf_beamline}{Mirror._reflectivity}

\api[\texttt{Filter}; \texttt{transfer}]{\texttt{Filter(material, density\_g\_cm3, thickness\_um, label="")}; \texttt{transfer(energies\_kev)}}
A slab (window, filter, air path) with Beer--Lambert transmission
$T(E) = \exp[-(\mu/\rho)(E)\,\rho\,t]$, using the total mass attenuation coefficient of the
compound (\fn{CS_Total_CP}: photoabsorption plus coherent and incoherent scattering, since any
interaction removes the photon from the beam). Energies without data are fully absorbed.
\pycode{xrf_beamline}{Filter}
\pycode{xrf_beamline}{Filter.__post_init__}
\pycode{xrf_beamline}{Filter.transfer}

\api[\texttt{Slit}; \texttt{transfer}]{\texttt{Slit(distance\_m, vertical\_gap\_mm=None, horizontal\_gap\_mm=None, label="")}; \texttt{transfer}}
A pair of slits: an aperture without spectral effect (\fn{transfer} returns ones). A gap of
\fn{None} is open.
\pycode{xrf_beamline}{Slit}
\pycode{xrf_beamline}{Slit.__post_init__}
\pycode{xrf_beamline}{Slit.transfer}

\subsection{Monochromatic sources and units}

\api[\texttt{tube\_lines}]{\texttt{tube\_lines(anode, lines=("KL3", "KL2"))}}
Energies (eV) and relative intensities of X-ray tube anode lines, by default the K$\alpha_1$/
K$\alpha_2$ doublet that a crystal or multilayer monochromator passes (\fn{tube_lines("Mo")}:
17\,479 and 17\,374\,eV). The intensities are the radiative rates (\fn{RadRate}), normalised
to sum 1; this is correct for lines that share the same initial vacancy, such as K$\alpha_1$
and K$\alpha_2$. The constant \fn{OpticalElement}, the union type of the optics, precedes it
in the file.
\begin{hyp}
The monochromator passes both lines with the same efficiency and suppresses K$\beta$ and the
bremsstrahlung completely.
\end{hyp}
\pycode{xrf_beamline}{tube_lines}

\api[\texttt{monochromatic\_spectrum}]{\texttt{monochromatic\_spectrum(energies\_ev, weights=(1.0,), width\_ev=5.0, step\_ev=1.0)}}
Builds a narrow incident spectrum from Gaussian lines of standard deviation \fn{width_ev},
for the functions that expect an incident spectrum (\fn{element_line_groups},
\fn{MatrixModel}). The density per eV, normalised so that each line has area $w_i$, is
converted to the per-0.1\,\% bandwidth convention, $N_{0.1\%} = 10^{-3}E\,\dd N/\dd E$.
\blk{Returns} \texttt{(energy\_kev, flux\_per\_0.1\%BW)} on a grid from $8\sigma$ below the lowest to
$8\sigma$ above the highest line.
\begin{hyp}
The width only has to be resolved by the grid; its value does not matter as long as the
cross-sections vary slowly across it (true away from absorption edges). With
\fn{MatrixModel}, pass \fn{sf_step_kev = step_ev/1000}: its secondary-fluorescence grid is
decimated to 40\,eV by default, which would miss a 5\,eV-wide line.
\end{hyp}
\pycode{xrf_beamline}{monochromatic_spectrum}

\api[\texttt{per\_unit\_energy}]{\texttt{per\_unit\_energy(energies\_kev, flux\_01bw)}}
Eq.~\eqref{eq:per-energy}: per-0.1\,\% bandwidth flux to a density per keV. Integrals over
$\dd E$ (flux-weighted cross-sections) must use this form.
\pycode{xrf_beamline}{per_unit_energy}

\subsection{Air}

\api[\texttt{AIR\_FORMULA}, \texttt{air\_density}]{\texttt{AIR\_FORMULA}, \texttt{air\_density(temperature\_c=20.0, pressure\_hpa=1013.25)}}
Dry air is written as an atom-count formula from its volume composition (N$_2$ 78.084\,\%,
O$_2$ 20.946\,\%, Ar 0.934\,\%, CO$_2$ 0.0417\,\%): $\mathrm{N}_{156.168}\mathrm{O}_{41.975}
\mathrm{Ar}_{0.934}\mathrm{C}_{0.0417}$ per 100 molecules, which reproduces the mass fractions
N 75.5\,\%, O 23.2\,\%, Ar 1.29\,\%. Its density follows the ideal-gas law,
\begin{equation}
  \rho = \frac{P\,M}{R\,T}, \qquad M = 28.9647\ \mathrm{g/mol},
\end{equation}
returned in g/cm$^3$ (1.204\,mg/cm$^3$ at 20\,\textdegree C and 1013.25\,hPa).
\begin{hyp}
Humidity is neglected. Moist air is slightly lighter than dry air, so at 50\,\% relative
humidity and 20\,\textdegree C this overestimates density and absorption by a few parts per
thousand \citep{Picard2008}. The pressure is the local value: about 990\,hPa at the ESRF,
2.5\,\% below the sea-level standard.
\end{hyp}
\pycode{xrf_beamline}{air_density}

\api[\texttt{air\_path}]{\texttt{air\_path(length\_m, temperature\_c=20.0, pressure\_hpa=1013.25)}}
A \fn{Filter} of air of the given length. A long air path is the strongest low-energy filter
of a white-beam line: 3\,m transmit about $10^{-7}$ at 5\,keV, 1.7\,\% at 8\,keV and 12\,\% at
10\,keV.
\pycode{xrf_beamline}{air_path}

\subsection{The assembled beamline}

\api[\texttt{Beamline}]{\texttt{Beamline(source, optics=[], sample\_distance\_m=None, vertical\_acceptance="auto", vertical\_offset\_urad=0.0, settings=\{\})}}
A source and a list of optical elements. The elements with a position act as apertures
(slits through their gaps, mirrors through their footprint); \fn{acceptance} maps each of them
back to an angle at the source, and the smallest vertical one sets the window of source
angles whose spectrum is used.
\begin{args}
\item[source] a \fn{BendingMagnetSource}.
\item[optics] mirrors, filters and slits. Positions come from \fn{distance_m}; the order of
  the list only sets the order of labels and transfers.
\item[sample\_distance\_m] source--sample distance (for \fn{trace} and \fn{describe}).
\item[vertical\_acceptance] \fn{"auto"} (from the apertures), \fn{"full"} (vertically
  integrated spectrum, Eq.~\eqref{eq:g1}) or a number (full acceptance in \textmu rad).
\item[vertical\_offset\_urad] centre of the accepted window relative to the orbit plane.
\item[settings] free-form provenance (motor positions read from a scan).
\end{args}
\pycode{xrf_beamline}{Beamline}

\api[\texttt{\_positioned}, \texttt{trace}]{\texttt{\_positioned()}, \texttt{trace(plane="vertical")}}
Paraxial ray tracing. A ray leaves the source on axis with unit angle; at each positioned
element its height $h$ (per radian of source angle, in m/rad) is accumulated, and every
mirror bent in \fn{plane} changes the ray angle by $-h/f$ (thin concave lens). Returns
\fn{[(element, h), ...]} in order of distance, plus \fn{(None, h_sample)} if the sample
distance is set. The height per unit source angle is what converts an aperture into an
angular acceptance at the source.
\pycode{xrf_beamline}{Beamline._positioned}
\pycode{xrf_beamline}{Beamline.trace}

\api[\texttt{acceptance}]{\texttt{acceptance()}}
For each plane and each aperture, the full angular acceptance at the source in \textmu rad,
$\Delta\psi = s/h$ with $s$ the gap or the mirror's $L\sin\theta$. Returns
\texttt{\{"vertical": \{limits, acceptance\_urad, limited\_by\}, "horizontal": \{...\}\}}, with
infinite acceptance when nothing limits a plane.
\pycode{xrf_beamline}{Beamline.acceptance}

\api[\texttt{effective\_vertical\_acceptance}]{\texttt{effective\_vertical\_acceptance()}}
The acceptance actually used for the spectrum: \fn{None} for \fn{"full"} or when no aperture
limits the beam, the computed value for \fn{"auto"}, or the number given.
\pycode{xrf_beamline}{Beamline.effective_vertical_acceptance}

\api[\texttt{resolved\_optics}]{\texttt{resolved\_optics()}}
Returns the optics with the grazing-angle spread of each focusing mirror filled in. An
elliptical mirror receives nearly parallel rays from the distant source and sends them
converging to the focus at distance $q$; across an illuminated beam of size $s$ the outgoing
direction varies by $s/q$, so the grazing angle varies by
\begin{equation}
  \Delta\theta = \pm\frac{s}{4q}
\end{equation}
\citep{Kirkpatrick1948}. $s$ is the accepted angle in the mirror's plane times the height per
angle at the mirror, capped by the mirror's own aperture. The reflectivity is then averaged
over this range, which smooths the high-energy cut-off.
\pycode{xrf_beamline}{Beamline.resolved_optics}

\api[\texttt{describe}]{\texttt{describe()}}
A dictionary summarising the geometry: vertical and horizontal acceptances and the aperture
limiting each, every aperture's limit, the KB angle spreads, the height per \textmu rad at
each element, the acceptance and offset used.
\pycode{xrf_beamline}{Beamline.describe}

\api[\texttt{spectrum}, \texttt{spectral\_density}, \texttt{band\_limits}]{\texttt{spectrum(energies\_kev)}, \texttt{spectral\_density(energies\_kev)}, \texttt{band\_limits(energies\_kev, frac=0.01)}}
\fn{spectrum} returns \fn{(raw, transmitted, transfers)}: the accepted source flux, the flux at
the sample (the product of every mirror and filter transfer), and the list of
\fn{(label, transfer)}, all per mrad and per 0.1\,\%\,BW. \fn{spectral_density} converts the
transmitted flux to photons/s/mrad/keV. \fn{band_limits} gives the energies where the
transmitted flux first and last exceeds \fn{frac} of its maximum---a compact statement of the
usable band (about 5--25\,keV on BM32).
\pycode{xrf_beamline}{Beamline.spectrum}
\pycode{xrf_beamline}{Beamline.spectral_density}
\pycode{xrf_beamline}{Beamline.band_limits}

\subsection{BM32}

\api[\texttt{BM32\_LAYOUT}, \texttt{BM32\_WHITE\_BEAM}]{\texttt{BM32\_LAYOUT}, \texttt{BM32\_WHITE\_BEAM(air\_path\_m=3.3, temperature\_c=20.0, pressure\_hpa=1013.25, current\_a=0.2, magnetic\_field\_t=0.86, *, vg1=0.3, hg1=0.2, vg2=None, hg2=None, vg3=0.2, hg3=0.2, m1\_angle\_mrad=3.062, m2\_angle\_mrad=2.551, m1\_radius\_km=71.4, m2\_radius\_km=145.7, kbv\_angle\_mrad=3.0, kbh\_angle\_mrad=3.0, kb\_angle\_spread=True, vertical\_acceptance="auto", vertical\_offset\_urad=0.0)}}
The BM32 white-beam model. \fn{BM32_LAYOUT} holds the distances from the source (m) of the
primary slits s1 (26.368), the bent Ir mirrors m1 (28.309) and m2 (31.732), the guard slits s2
(35.370), the slits s3 that define the KB illumination (44.485), the KB mirrors (vertical
44.645, horizontal 44.880) and the sample (45.0), and the mirror lengths. The function builds:
s1; m1 and m2 (Ir, 22.56\,g/cm$^3$, bent); two 500\,\textmu m Be windows; s2; s3; KB-V and KB-H
(Ir, focusing, with angle spread if \fn{kb_angle_spread}); and the incident air path.
\begin{args}
\item[air\_path\_m] incident air path (3.3\,m measured on BM32); 0 removes it. Not to be
  confused with the sample--detector path of the detector model.
\item[vg1 \dots\ hg3] slit gaps (mm), motors of the same names; \fn{None} = open.
\item[m1\_angle\_mrad, m2\_angle\_mrad, m1\_radius\_km, m2\_radius\_km] mirror angles
  (motors \fn{ma1}, \fn{ma2}) and bending radii (\fn{mc1}, \fn{mc2}).
\item[kbv\_angle\_mrad, kbh\_angle\_mrad] KB grazing angles, not recorded in the scan files.
\item[others] as in \fn{Beamline} and \fn{BendingMagnetSource}.
\end{args}
With the default slits the vertical acceptance is a few \textmu rad, set by s3: the sample sees
the on-axis spectrum, which raised the Ga/As sensitivity ratio by 4.5\,\% on the 2026 maps
compared with the vertically integrated spectrum.
\pycode{xrf_beamline}{BM32_WHITE_BEAM}

\api[\texttt{bm32\_from\_h5}]{\texttt{bm32\_from\_h5(h5\_path, scan=None, **overrides)}}
The BM32 beamline with the settings recorded in a BLISS scan: slit gaps
(\fn{vg1 hg1 vg2 hg2 vg3 hg3}), mirror angles (\fn{ma1 ma2}) and radii (\fn{mc1 mc2}) from
\fn{<scan>/instrument/positioners}, and the ring current from
\fn{<scan>/instrument/machine/current} (mA, converted to A). Missing or non-numeric values
(BLISS writes \fn{ERR} for disconnected motors) keep their defaults with a warning; keyword
arguments override anything. The source file and scan are stored in \fn{settings}.
\pycode{xrf_beamline}{bm32_from_h5}
